\begin{thebibliography}{99}
\small

\bibitem[{Arthur et~al.(1998)}]{arthur1998}
Arthur, W., Bennett, W., Stanush, P.~L. y McNelly, T.~L. (1998).
\newblock Factors that influence skill decay and retention: A quantitative review and analysis.
\newblock \emph{Human Performance}, 11(1):57--101.

\bibitem[{Barzagar~Nazari y Ebersbach(2019)}]{barzagar2019}
Barzagar~Nazari, K. y Ebersbach, M. (2019).
\newblock Distributing mathematical practice of third and seventh graders: Applicability of the spacing effect in the classroom.
\newblock \emph{Applied Cognitive Psychology}, 33(2):288--298.

\bibitem[{Cen et~al.(2006)}]{cen2006}
Cen, H., Koedinger, K. y Junker, B. (2006).
\newblock Learning factors analysis: A general method for cognitive model evaluation and improvement.
\newblock En \emph{Intelligent Tutoring Systems (ITS 2006)}, pp.~164--175. Springer.

\bibitem[{Cepeda et~al.(2006)}]{cepeda2006}
Cepeda, N.~J., Pashler, H., Vul, E., Wixted, J.~T. y Rohrer, D. (2006).
\newblock Distributed practice in verbal recall tasks: A review and quantitative synthesis.
\newblock \emph{Psychological Bulletin}, 132(3):354--380.

\bibitem[{Cepeda et~al.(2008)}]{cepeda2008}
Cepeda, N.~J., Vul, E., Rohrer, D., Wixted, J.~T. y Pashler, H. (2008).
\newblock Spacing effects in learning: A temporal ridgeline of optimal retention.
\newblock \emph{Psychological Science}, 19(11):1095--1102.

\bibitem[{Choffin et~al.(2019)}]{choffin2019}
Choffin, B., Popineau, F., Bourda, Y. y Vie, J.-J. (2019).
\newblock DAS3H: Modeling student learning and forgetting for optimally scheduling distributed practice of skills.
\newblock En \emph{Proceedings of the 12th International Conference on Educational Data Mining (EDM 2019)}. arXiv:1905.06873.

\bibitem[{Choffin et~al.(2021)}]{choffin2021}
Choffin, B., Popineau, F. y Bourda, Y. (2021).
\newblock Extending adaptive spacing heuristics to multi-skill items.
\newblock \emph{Journal of Educational Data Mining}, 13(3).

\bibitem[{Corbett y Anderson(1994)}]{corbett1994}
Corbett, A.~T. y Anderson, J.~R. (1994).
\newblock Knowledge tracing: Modeling the acquisition of procedural knowledge.
\newblock \emph{User Modeling and User-Adapted Interaction}, 4(4):253--278.

\bibitem[{de~la Torre(2011)}]{delatorre2011}
de~la Torre, J. (2011).
\newblock The generalized DINA model framework.
\newblock \emph{Psychometrika}, 76(2):179--199.

\bibitem[{Doignon y Falmagne(1985)}]{doignon1985}
Doignon, J.-P. y Falmagne, J.-C. (1985).
\newblock Spaces for the assessment of knowledge.
\newblock \emph{International Journal of Man-Machine Studies}, 23(2):175--196.

\bibitem[{Hunziker et~al.(2019)}]{hunziker2019}
Hunziker, A., Chen, Y., Mac~Aodha, O., Gomez-Rodriguez, M., Krause, A., Perona, P., Yue, Y. y Singla, A. (2019).
\newblock Teaching multiple concepts to a forgetful learner.
\newblock En \emph{Advances in Neural Information Processing Systems 32 (NeurIPS 2019)}.

\bibitem[{Junker y Sijtsma(2001)}]{junker2001}
Junker, B.~W. y Sijtsma, K. (2001).
\newblock Cognitive assessment models with few assumptions, and connections with nonparametric item response theory.
\newblock \emph{Applied Psychological Measurement}, 25(3):258--272.

\bibitem[{Koedinger et~al.(2012)}]{koedinger2012}
Koedinger, K.~R., Corbett, A.~T. y Perfetti, C. (2012).
\newblock The Knowledge-Learning-Instruction framework: Bridging the science-practice chasm to enhance robust student learning.
\newblock \emph{Cognitive Science}, 36(5):757--798.

\bibitem[{Lindsey et~al.(2014)}]{lindsey2014}
Lindsey, R.~V., Shroyer, J.~D., Pashler, H. y Mozer, M.~C. (2014).
\newblock Improving students' long-term knowledge retention through personalized review.
\newblock \emph{Psychological Science}, 25(3):639--647.

\bibitem[{Nemhauser et~al.(1978)}]{nemhauser1978}
Nemhauser, G.~L., Wolsey, L.~A. y Fisher, M.~L. (1978).
\newblock An analysis of approximations for maximizing submodular set functions---I.
\newblock \emph{Mathematical Programming}, 14(1):265--294.

\bibitem[{Pavlik et~al.(2009)}]{pavlik2009}
Pavlik, P.~I., Cen, H. y Koedinger, K.~R. (2009).
\newblock Performance factors analysis: A new alternative to knowledge tracing.
\newblock En \emph{Artificial Intelligence in Education (AIED 2009)}, pp.~531--538. IOS Press.

\bibitem[{Pyc y Rawson(2009)}]{pyc2009}
Pyc, M.~A. y Rawson, K.~A. (2009).
\newblock Testing the retrieval effort hypothesis: Does greater difficulty correctly recalling information lead to higher levels of memory?
\newblock \emph{Journal of Memory and Language}, 60(4):437--447.

\bibitem[{Reddy et~al.(2016)}]{reddy2016}
Reddy, S., Labutov, I., Banerjee, S. y Joachims, T. (2016).
\newblock Unbounded human learning: Optimal scheduling for spaced repetition.
\newblock En \emph{Proceedings of the 22nd ACM SIGKDD International Conference on Knowledge Discovery and Data Mining}, pp.~1815--1824.

\bibitem[{Rohrer et~al.(2015)}]{rohrer2015}
Rohrer, D., Dedrick, R.~F. y Stershic, S. (2015).
\newblock Interleaved practice improves mathematics learning.
\newblock \emph{Journal of Educational Psychology}, 107(3):900--908.

\bibitem[{Rohrer y Taylor(2007)}]{rohrer2007}
Rohrer, D. y Taylor, K. (2007).
\newblock The shuffling of mathematics problems improves learning.
\newblock \emph{Instructional Science}, 35(6):481--498.

\bibitem[{Settles y Meeder(2016)}]{settles2016}
Settles, B. y Meeder, B. (2016).
\newblock A trainable spaced repetition model for language learning.
\newblock En \emph{Proceedings of the 54th Annual Meeting of the Association for Computational Linguistics (ACL 2016)}, pp.~1848--1858.

\bibitem[{Skycak(2023)}]{skycak2023}
Skycak, J. (2023).
\newblock Individualized spaced repetition in hierarchical knowledge structures.
\newblock \url{https://www.justinmath.com/individualized-spaced-repetition-in-hierarchical-knowledge-structures/}.

\bibitem[{Su et~al.(2023)}]{su2023}
Su, J., Ye, J., Nie, L., Cao, Y. y Chen, Y. (2023).
\newblock Optimizing spaced repetition schedule by capturing the dynamics of memory.
\newblock \emph{IEEE Transactions on Knowledge and Data Engineering}, 35(10):10085--10097.

\bibitem[{Tabibian et~al.(2019)}]{tabibian2019}
Tabibian, B., Upadhyay, U., De, A., Zarezade, A., Schölkopf, B. y Gomez-Rodriguez, M. (2019).
\newblock Enhancing human learning via spaced repetition optimization.
\newblock \emph{Proceedings of the National Academy of Sciences}, 116(10):3988--3993.

\bibitem[{Tatsuoka(1983)}]{tatsuoka1983}
Tatsuoka, K.~K. (1983).
\newblock Rule space: An approach for dealing with misconceptions based on item response theory.
\newblock \emph{Journal of Educational Measurement}, 20(4):345--354.

\bibitem[{Upadhyay et~al.(2021)}]{upadhyay2021}
Upadhyay, U., Lancashire, G., Moser, C. y Gomez-Rodriguez, M. (2021).
\newblock Large-scale randomized experiments reveals that machine learning-based instruction helps people memorize more effectively.
\newblock \emph{npj Science of Learning}, 6:26.

\bibitem[{Woźniak(1990)}]{wozniak1990}
Woźniak, P.~A. (1990).
\newblock \emph{Optimization of learning}.
\newblock Tesis de maestría, Universidad Tecnológica de Poznań.

\bibitem[{Ye et~al.(2022)}]{ye2022}
Ye, J., Su, J. y Cao, Y. (2022).
\newblock A stochastic shortest path algorithm for optimizing spaced repetition scheduling.
\newblock En \emph{Proceedings of the 28th ACM SIGKDD Conference on Knowledge Discovery and Data Mining}, pp.~4381--4390.

\end{thebibliography}
